% 第10章 另类数据与微观结构
\section{另类数据与微观结构：信号的上游}\label{sec:altdata}

前九章讨论的方法，都需要“特征”。这一章讨论特征从哪来。另类数据（alternative data）和微观结构数据是这一代量化最重要的信息增量来源——它们不是新算法，而是\textbf{把原本没有被数字化、或者没有被充分利用的信息，变成可计算的特征}。用第 \ref{sec:foundations} 章的话说，它们动的是“信息存量的天花板”，而不是“把已有信息利用得更充分”。

\subsection{另类数据的谱系}

\begin{itemize}
  \item \textbf{卫星与遥感}：停车场车流量、油罐储量、夜间灯光、农作物长势。Katona 等（2018）\cite{katona2018} 讨论了卫星数据进入资本市场的机制与后果。
  \item \textbf{消费与支付}：信用卡交易、App 使用、电商价格与库存、外卖订单。
  \item \textbf{人力与组织}：招聘职位数、员工流动、专利与研发活动。Zhu（2019）\cite{zhu2019} 研究了大数据（含招聘与员工评价数据）如何影响公司治理。
  \item \textbf{众包与网络}：众包盈利预测、社交媒体、搜索引擎热度。Jame 等（2016）\cite{jame2016} 检验了众包盈利预测的预测力。
  \item \textbf{资金与持仓流}：ETF 申赎、机构持仓、期权暗池流。
\end{itemize}

\subsection{微观结构：最接近“物理真实”的数据}

如果要在这份综述里挑一类数据，说它“最适合与机器学习结合且最可信”，会是限价订单簿与逐笔成交。原因有三：样本量以十亿计；生成机制相对稳定（撮合规则是明确的）；且预测目标与执行决策直接挂钩。DeepLOB\cite{zhang2019deeplob}、Tsantekidis 等（2017）\cite{tsantekidis2017} 一类工作证明了深度学习能从订单簿中提取出可用的短时价格信息；Ntakaris 等（2018）\cite{ntakaris2018fi2010} 提供了公开基准数据集，使得这一方向的复现性在六族之外一枝独秀。订单簿数据的“贵”在于获取与存储成本，而非稀缺——交易所逐笔数据可以买，只是很贵。

\subsection{落地要素：另类数据的四个隐性成本}

另类数据看起来是“多一类特征”这么简单，实际落地要跨越四道关：

\begin{enumerate}
  \item \textbf{许可与合规}。数据的使用权（尤其涉及个人信息的信用卡、定位数据）有严格的法律边界，很多数据源在合规审查下不可用。
  \item \textbf{覆盖率与偏差}。一种支付数据可能只覆盖某一人群或某一地区，用它预测全体公司，会引入样本选择偏差。
  \item \textbf{时点与修订}。数据源本身可能事后修订历史值，若不保存“当时看到的那份”，就会产生前视偏差。
  \item \textbf{成本与容量}。另类数据订阅费用高昂，且信号容量（能承载的资金规模）有限，必须与成本一起评估。
\end{enumerate}

\subsection{制度约束的特殊性：以 A 股为例}\label{sec:china}

前文的讨论以美股语境为主，但交易制度会把同一套方法的结果改得面目全非。中文市场的读者尤其需要注意下面几点差异。

\begin{itemize}
  \item \textbf{涨跌停}：主板 $\pm10\%$，创业板与科创板 $\pm20\%$。触及涨跌停的股票当日不可买入（涨停）或不可卖出（跌停）；在极端行情中大量个股同时触板，组合调仓会在最需要它的时刻被冻结。若回测忽略这一点，等于假设自己在最拥挤的时刻仍能自由成交。
  \item \textbf{T+1}：当日买入的股票次日才能卖出。任何依赖日内反转或快速止盈的短周期策略，在 A 股必须重写持有期约束。
  \item \textbf{融券约束}：个股做空渠道有限、成本高、券源不稳。教科书式的“多空对称”组合在 A 股往往不可实现，多数策略实际只能做多头增强或以股指期货对冲。
  \item \textbf{散户占比与情绪}：情绪与舆情类信号在 A 股的边际信息可能更强，但衰减也更快；这既解释了部分信号的有效性，也提示其易拥挤。
  \item \textbf{数据可得性}：本地通达信日线（含主力连续与持仓量）可作为低成本的历史底座；财务数据必须按实际披露日对齐——A 股披露集中在特定窗口，容易在截面上造成“同一时点信息高度相关”，进而高估模型的独立信息量。
\end{itemize}

\textbf{复现上的含义很直接}：把在美股验证过的方法搬到 A 股，必须重做制度掩码与成本口径（对应第 \ref{sec:repro_protocol} 章 P0 的“交易制度掩码”条款）；不重做，结论之间不可比，也不可用。

\subsection{复现状态与失效边界}

\textbf{另类数据是六族之外可复现性最差的一类，原因与算法无关：数据本身是私有的。}私有数据上的结果，外部研究者无法独立验证，只能依赖发表方的自述。这不是学术缺陷，而是这类研究的天然属性，但它意味着\textbf{任何基于私有另类数据的高收益声明，都应被视为“无法独立验证的经验”而非“科学结论”}。

失效边界里最该警惕的是\textbf{衰减速度}：另类数据一旦被证明有效，就会有更多资金涌入（Grossman--Stiglitz 的老故事），超额收益迅速收窄。卫星图像、信用卡数据在 2010 年代中后期曾是显著 alpha，如今多数已相当拥挤。相反，微观结构数据因为需要工程能力与基础设施，衰减相对慢，这也是它长期有价值的原因。本章要记住的一句话是：\textbf{另类数据卖的不是数据，是时间差——它值钱只因为你比别人早用上，而一旦人人都在用，它就变成了新的“公开信息”并被价格吞掉。}
